\chapter{2015年重庆卷（文）}

\section{单选题}

\begin{problem}
已知集合 \(A = \{1, 2, 3\}\)，\(B = \{1, 3\}\)，则 \(A \cap B =\)（\quad）

\choices
    {{2}}
    {\( \{1,2\} \)}
    {\(\{1,3\}\)}
    {\(\{1,2,3\}\)}
\end{problem}

\begin{problem}
“ \( z = 1 \) ”是“ \( x^{2} - 2x + 1 = 0 \) ”的（\quad）

\choices
    {充要条件}
    {充分而不必要条件}
    {必要而不充分条件}
    {既不充分也不必要条件}
\end{problem}

\begin{problem}
函数 \( f(x) = \log_2 (x^2 + 2x - 3) \) 的定义域是（\quad）

\choices
    {\([-3, 1]\)}
    {\((-3, 1)\)}
    {\((-\infty, -3] \cup \{1, +\infty\}\)}
    {\((-\infty, -3) \cup (1, +\infty)\)}
\end{problem}

\begin{problem}
重庆市 2013 年各月的平均气温(\(^{\circ}\mathrm{C}\)) 数据的茎叶图如下:

则这组数据的中位数是（\quad）

\choices
    {19}
    {20}
    {21.5}
    {23}
\end{problem}

\begin{problem}
某几何体的三视图如图所示，则该几何体的体积为（\quad）

正视图

左视图

俯视图（\quad）

\choices
    {\(\frac{1}{3} + 2\pi\)}
    {\(\frac{13\pi}{6}\)}
    {\(\frac{7\pi}{3}\)}
    {\(\frac{5\pi}{2}\)}
\end{problem}

\begin{problem}
若 \(\tan \alpha = \frac{1}{3}, \tan (\alpha + \beta) = \frac{1}{2}\)，则 \(\tan \beta =\)（\quad）

\choices
    {\( \frac{1}{7} \)}
    {\( \frac{1}{6} \)}
    {\( \frac{5}{7} \)}
    {\( \frac{5}{6} \)}
\end{problem}

\begin{problem}
已知非零向量 \(a, b\) 满足 \(|b| = 4|a|\)，且 \(a \perp (2a + b)\)，则 \(a\) 与 \(b\) 的夹角为（\quad）

\choices
    {\( \frac{\pi}{3} \)}
    {\( \frac{\pi}{2} \)}
    {\( \frac{2\pi}{3} \)}
    {\(\frac{5\pi}{6}\)}
\end{problem}

\begin{problem}
执行如图所示的程序框图，则输出 s 的值为（\quad）

\choices
    {\(\frac{3}{4}\)}
    {\(\frac{5}{6}\)}
    {\(\frac{11}{12}\)}
    {\(\frac{25}{24}\)}
\end{problem}

\begin{problem}
设双曲线 \(\frac{x^2}{a^2} - \frac{y^2}{b^2} = 1 (a > 0, b > 0)\) 的右焦点是 \(F\)、左、右顶点分别是 \(A_1, A_2\)，设 \(F\) 作 \(A_1A_2\) 的垂线与双曲线交于 \(B\)、\(C\) 两点. 若 \(A_1B \perp A_2C\)，则该双曲线的渐近线斜率为（\quad）

\choices
    {\(\pm \frac{1}{2}\)}
    {\(\pm \frac{\sqrt{2}}{2}\)}
    {\(\pm 1\)}
    {\(\pm \sqrt{2}\)}
\end{problem}

\begin{problem}
若不等式组 \( \left\{\begin{aligned}x+y-2&\leqslant0,\\ x+2y-2&\geqslant0,\end{aligned}\right. \) 表示的平面区域为三角形，且其面积等于 \( \frac{4}{3} \) ，则 m 的值为（\quad）

\choices
    {\(-3\)}
    {1}
    {\( \frac{4}{3} \)}
    {3}
\end{problem}

\section{填空题}

\begin{problem}
复数 \((1 + 2\i)\) 的实部为
\end{problem}

\begin{problem}
若点 \(P(1,2)\) 在以坐标原点为圆心的圆上，则该圆在点 \(P\) 处的切线方程为 \fillinblank{} \fillinblank{}.
\end{problem}

\begin{problem}
设 \(\triangle ABC\) 的内角 \(A, B, C\) 的对边分别为 \(a, b, c\)，且 \(a = 2, \cos C = -\frac{1}{4}\)，\(3\sin A = 2\sin B\)，则 \(c = \fillinblank{}\).
\end{problem}

\begin{problem}
设 \(a, b > 0, a + b = 5\)，则 \(\sqrt{a + 1} + \sqrt{b + 3}\) 的最大值为 \fillinblank{} \fillinblank{}.
\end{problem}

\begin{problem}
在区间 [0, 5] 上随机地选择一个数 \( p \)，则方程 \( x^2 + 2px + 3p - 2 = 0 \) 有两个负根的概率为 \fillinblank{}.
\end{problem}

\section{解答题}
\begin{problem}
已知等差数列 \(\{a_{n}\}\) 满足 \(a_3 = 2\)，前3项和 \(S_{3} = \frac{9}{2}\). (1) 求 \( \{a_{n}\} \) 的通项公式; (2) 设等比数列 \( \{b_{n}\} \) 满足 \( b_{1}=a_{1}, b_{4}=a_{15} \) ，求 \( \{b_{n}\} \) 的前 n 项和 \( T_{n} \)  \fillinblank{}.
\end{problem}

\begin{problem}
随着我国经济的发展，居民的储蓄存款逐年增长. 设某地区城乡居民人民币储蓄存款 (年底余额) 如下表: 年份 2010 2011 2012 2013 2014 时间代号t 1 2 3 4 5 储蓄存款y(千亿元) 5 6 7 8 10 (1) 求 y 关于 t 的回归方程 \( \hat{y} = bt + \hat{a} \) ; (2) 用所求回归方程预测该地区 2015 年 (=6) 的人民币储蓄存款. 附: 回归方程 \( \hat{y} = \hat{b} t + \hat{a} \) 中，\( \hat{b} = \frac{\sum_{i=1}^{n} t_{i} y_{i} - n \overline{t} \overline{y}}{\sum_{i=1}^{n} t_{i}^{2} - n \overline{t}^{2}} \) ，\( \hat{a} = \overline{y} - \hat{b} \overline{t} \)  \fillinblank{}.
\end{problem}

\begin{problem}
已知函数 \( f(x) = \frac{1}{2}\sin 2x - \sqrt{3}\cos^2 x \). (1) 求 \( f(x) \) 的最小正周期和最小值. (2) 将函数 \( f(x) \) 的图象上每一点的横坐标伸长到原来的两倍，纵坐标不变，得到函数 \( g(x) \) 的图象，当 \( x \in \left[\frac{\pi}{2}, \pi\right] \) 时，求 \( g(x) \) 的值域 \fillinblank{}.
\end{problem}

\begin{problem}
如图，三棱锥 \(P - ABC\) 中，平面 \(PAC \perp\) 平面 \(ABC\)，\(\angle ABC = \frac{\pi}{3}\)，点 \(D\)，\(E\) 在线段 \(AC\) 上，且 \(AD = DE = EC = 2\)，\(PD = PC = 4\)，点 \(\overrightarrow{P}\) 在线段 \(AB\) 上，且 \(EF // BC\). (1) 证明: \(AB \perp\) 平面 \(PFE\); (2) 若四棱锥 P-DFBC 的体积为 7，求线段 BC 的长 \fillinblank{}.
\end{problem}

\begin{problem}
如图，椭圆 \(\frac{x^2}{a^2} + \frac{y^2}{b^2} = 1 (a > b > 0)\) 的左、右焦点分别为 \(F_1, F_2\)，过 \(F_2\) 的直线交椭圆于 \(P, Q\) 两点，且 \(PQ \perp PF_1\). (1) 若 \( |PF_1| = 2 + \sqrt{2} \)，\( |PF_2| = 2 - \sqrt{2} \)，求椭圆标准方程; (2) 若 \( |PQ| = \lambda |PF_1| \)，且 \( \frac{9}{4} \leqslant \lambda < \frac{4}{3} \)，试确定椭圆离心率 \( e \) 的取值范围 \fillinblank{}.
\end{problem}

\begin{problem}
已知函数 \( f(x) = ax^2 + x^2 (a \in \R) \) 在 \( x = -\frac{4}{3} \) 处取得极值. (1) 确定 a 的值; (2) 若 \( g(x) = f(x) \e^{x} \) ，讨论 \( g(x) \) 的单调性 \fillinblank{}.
\end{problem}

